\documentclass[11pt]{article}
\usepackage[margin=1.1in]{geometry}
\usepackage{amsmath}
\usepackage{graphicx}
\usepackage{booktabs}
\usepackage{enumitem}
\usepackage[hidelinks]{hyperref}

% ============================================================================
% REPORT SKELETON — Measuring and Removing an Object's Light Footprint
% How to use: every section has a % PLAN comment telling you what to write and
% which material to pull from. Write in your own words; the proposal supplies
% Sections 1-3 almost verbatim. Target length: 8-12 pages including figures.
% ============================================================================

\title{Measuring and Removing an Object's Light Footprint:\\
\large Reflection-Aware Object Removal in 3D Gaussian Splatting via Render--Edit--Refit}
\author{Yonghao Lee\\[2pt]
{\normalsize Deep Learning for 3D Computer Vision --- Final Report}\\
{\normalsize Hebrew University of Jerusalem \quad $\cdot$ \quad Instructor: Sagie Benaim}}
\date{September 2026}

\begin{document}
\maketitle

\begin{abstract}
% PLAN: 6-8 sentences. (1) problem: deleting an object from 3DGS leaves its
% shadow/reflection/bleed; (2) approach: render-edit-refit with a video prior;
% (3) benchmark: synthetic scene with exact clean plates + footprint regions;
% (4) key negative finding: generalist prior substitutes, even reconstructs the
% object from its footprint; (5) key positive: removal-specialized prior (ROSE)
% + 4-click SAM2 masks: footprint PSNR 14.2 -> 28.4 dB over deletion;
% (6) edits survive lifting to 3D (31.3 dB on held-out views).
TODO
\end{abstract}

\section{Introduction}
% PLAN: your proposal's Problem Definition, lightly extended. Vase picture
% first, L = L_direct + L_obj, goal sentence. Then ONE new paragraph
% summarizing contributions as a list of 4:
%   (1) a controlled benchmark with exact clean plates and footprint regions;
%   (2) a render-edit-refit pipeline requiring no COLMAP and no training;
%   (3) an empirical study of video priors as removers (substitution finding);
%   (4) a practical 4-click protocol matching oracle localization.
TODO

\section{Related Work}
% PLAN: copy your final proposal's five-family Related Work (removal/inpainting,
% footprint-aware, measuring-residuals, diffusion 3D editing, video priors) and
% update ONE thing: ROSE moves from "closest work" to "the editor we adopt";
% add the sentence: "ROSE removes side effects within a single 2D video; we
% study whether such a prior can remove them from a 3D scene, and measure per
% effect what survives the lift."
TODO

\section{Benchmark}
% PLAN: describe the Blender scene AND justify each design choice (this is
% where the war stories earn marks):
%  - polished dielectric floor: a metallic mirror floor shows no shadows
%    (no diffuse response) -- discovered via the eyeball check;
%  - separate vertical mirror panel: dielectric Fresnel reflection too weak
%    at this viewing angle for a headline effect;
%  - fill light: a mirror faithfully shows the UNLIT side of the object;
%  - with/without renders share cameras -> exact clean plates; differencing
%    outside the silhouette -> per-region footprint masks, no annotation;
%  - 81 frames at 832x480 chosen to match video-model native size (and 16n+1);
%  - every 8th frame held out: metrics measure novel-view quality in EVERY
%    experiment (interpolation, not extrapolation -- state this honestly).
% FIGURE 1: with-frame vs clean plate side by side (mid-arc).
TODO

\section{Method: Render--Edit--Refit}
% PLAN: your proposal's Approach section, in past tense, with the numbered
% 4-step pipeline. Add implementation facts: splatfacto 15k iters, random init
% (justify: exact poses + bounded synthetic scene; cite original 3DGS finding),
% masks white=edit, ROSE 50 steps, refit = identical training recipe on edited
% frames with unchanged poses. State the mask protocols precisely:
% object (SAM silhouette, dilated 25px), oracle (GT differencing -- labeled as
% upper bound), practical (4 clicks on one frame, SAM2 propagation, union,
% dilated 15px).
TODO

\section{Experiments}

\subsection{Reconstruction quality (ceiling)}
% PLAN: 45.7 dB / 0.992 SSIM / 0.021 LPIPS on held-out views; one sentence on
% why this matters (any later degradation is attributable to removal, not
% reconstruction). FIGURE 2: render vs GT side by side.
TODO

\subsection{The problem, quantified (M1: deletion)}
% PLAN: mask-lifting selection validated against ground-truth geometry
% (2,888 selected, 100% precision, 90% recall; 318 interior blobs found only
% geometrically). Deletion result: sphere gone, mirror image INTACT.
% FIGURE 3: the ghost figure (deleted render vs clean plate).
TODO

\subsection{Main results}
\begin{table}[h]
\centering
\begin{tabular}{lcc}
\toprule
Method & PSNR full & PSNR footprint \\
\midrule
Reconstruction ceiling            & 45.7 & --   \\
M1: Gaussian deletion             & 23.1 & 14.2 \\
M3: VACE, object mask             & 21.6 & 14.2 \\
M3: VACE, oracle mask             & 18.2 & 13.1 \\
M3: ROSE, object mask             & 23.2 & 15.0 \\
M3: ROSE, oracle mask             & 31.1 & 28.4 \\
M3: ROSE, SAM2 masks (4 clicks)   & \textbf{31.3} & \textbf{28.4} \\
\bottomrule
\end{tabular}
\caption{PSNR (dB) vs.\ exact clean plates on held-out novel views, full-frame
and restricted to ground-truth footprint regions. All rows measured after
re-fitting a fresh 3DGS on the edited frames.}
\end{table}
% PLAN: walk the table row by row (you already know the story of each).
TODO

\subsection{Why the generalist prior fails: substitution, not removal}
% PLAN: the VACE story with FIGURE 4 (gray sphere / cyan sphere / orange
% triptych). Key sentences:
%  - object-mask run REBUILT the sphere from its shadow+reflection: direct
%    visual evidence the footprint encodes the object (connect to Remove360's
%    semantic residuals);
%  - oracle run substituted coherently (cyan reflection matched!) -- the prior
%    understands object-reflection coupling but "generate emptiness" is not an
%    operation it learned;
%  - prompting cannot override the mask's shape prior (3 prompts, 3 objects).
TODO

\subsection{Multi-view consistency of independent edits}
% PLAN: the extension resolved by measurement: ROSE edits 81 frames
% independently, yet the refit reaches 31.3 dB on held-out views -> the 3D
% representation absorbs residual disagreement at this operating point;
% iterative dataset update and warped noise were therefore not needed.
% (One honest sentence: at lower edit quality they may still matter.)
TODO

\subsection{The practical protocol}
% PLAN: 4 clicks on one frame -> SAM2 propagates 4 regions across 81 frames
% (note: shadows and floor sheen are outside SAM2's comfort zone; minor drift
% observed, absorbed by mask union + dilation; wall bleed was too faint to
% click yet did not measurably hurt). Result ties the oracle.
% FIGURE 5: annotation frame with 4 X marks + one propagated-mask overlay.
% FIGURE 6: original vs ROSE+SAM2 result (the empty room).
TODO

\section{Limitations and Lessons}
% PLAN: honest list, prose not bullets if you prefer:
% single synthetic scene, one object; 120-degree arc (interpolation only);
% no keep-loss in the refit (unedited-region drift bounded by full-frame PSNR);
% oracle masks use GT (but practical protocol matched them); metrics are
% footprint-combined, not yet per-effect (shadow vs reflection vs bleed) --
% the masks to split them exist, listed as future work.
% Optionally 2-3 sentences of engineering lessons (version pinning; verify
% which inputs reached the code before theorizing).
TODO

\section{Conclusion}
% PLAN: 4-5 sentences. The headline sentence with numbers; the substitution
% finding as the second takeaway; benchmark released with scripts (link repo).
TODO

\section*{Contribution Statement}
% Required by course if pairs; for solo, one sentence + acknowledge tools:
% "All experiments and writing by the author. Coding assistance and debugging
% support from an AI assistant (Claude); all results verified by the author."
TODO

\section*{References}
% PLAN: copy from the final proposal (20 entries), add SAM 2 if not present.
TODO

\end{document}
